\documentclass[a4paper,reqno,10pt]{amsart}

\usepackage[
  colorlinks=true,
  linkcolor=blue,
  citecolor=blue,
  urlcolor=blue
]{hyperref}
\usepackage{a4wide}
\usepackage{amssymb,amsmath,amsthm,amsfonts,mathtools}
\usepackage{enumitem}
\usepackage[all]{xy}
\setlist[enumerate,1]{label={\upshape(\arabic*)}}
\setlist[enumerate,2]{label={\upshape(\alph*)}}

\newtheorem{theorem}{Theorem}[section]
\newtheorem{proposition}[theorem]{Proposition}
\newtheorem{lemma}[theorem]{Lemma}
\newtheorem{corollary}[theorem]{Corollary}
\newtheorem{conjecture}[theorem]{Conjecture}
\newtheorem*{theoremA}{Theorem A}
\newtheorem*{corollaryB}{Corollary B}
\theoremstyle{definition}
\newtheorem{definition}[theorem]{Definition}
\newtheorem*{organization}{Organization}
\newtheorem*{conventions}{Conventions and notation}
\newtheorem*{aiuse}{Use of AI}
\numberwithin{equation}{section}

\newcommand{\Hom}{\operatorname{Hom}}
\newcommand{\Ext}{\operatorname{Ext}}
\newcommand{\Tor}{\operatorname{Tor}}
\newcommand{\pd}{\operatorname{pd}}
\newcommand{\rad}{\operatorname{rad}}
\newcommand{\modu}{\operatorname{mod}}
\newcommand{\ARC}{\textup{(ARC)}}
\newcommand{\GNC}{\textup{(GNC)}}
\newcommand{\AGC}{\textup{(AGC)}}
\newcommand{\NC}{\textup{(NC)}}
\newcommand{\TCone}{\textup{(TC1)}}
\newcommand{\TCtwo}{\textup{(TC2)}}
\newcommand{\GPC}{\textup{(GPC)}}
\renewcommand{\baselinestretch}{1}

\title[Tachikawa's second conjecture and the Auslander--Reiten conjecture]
{Tachikawa's second conjecture implies\\
the Auslander--Reiten conjecture}

\author[H.\ Enomoto]{Haruhisa Enomoto}
\address{Parakeet Inc., Japan}
\email{haruhisa.enomoto.math@gmail.com}

\subjclass[2020]{16E30, 16D50, 16G10}
\keywords{Auslander--Reiten conjecture, Tachikawa conjecture,
Nakayama conjecture, artin algebra, self-injective algebra, trivial extension}
\date{\today}

\begin{document}

\begin{abstract}
We prove that Tachikawa's second conjecture implies the Auslander--Reiten
conjecture for artin algebras over a commutative artinian ring.  The proof
uses the two-fold trivial extension of an algebra.  Together with known
implications, it follows that the Auslander--Reiten conjecture, the
generalized Nakayama conjecture, the Auslander--Gorenstein conjecture, the
Nakayama conjecture, the Gorenstein-projective conjecture, and Tachikawa's
second conjecture are all equivalent, and that Tachikawa's second conjecture
implies his first.
\end{abstract}

\maketitle

\section{Introduction}

Auslander and Reiten proposed the following conjecture in their study of the
generalized Nakayama conjecture \cite{AR}.

\begin{conjecture}[Auslander--Reiten {\cite{AR}}]
\label{conj:arc}
Let $\Lambda$ be an artin algebra, and let $X$ be a finitely generated right
$\Lambda$-module.  If
\begin{equation}\label{eq:arc}
 \Ext_\Lambda^i(X,X\oplus\Lambda)=0\qquad\text{for every }i>0,
\end{equation}
then $X$ is projective.
\end{conjecture}

If an artin algebra $\Lambda$ is self-injective, then $\Lambda_\Lambda$ is
injective, so
$\Ext_\Lambda^i(X,\Lambda)=0$ for every $X$ and every $i>0$.  For
self-injective algebras, Conjecture~\ref{conj:arc} therefore becomes the
following conjecture of Tachikawa, known as his second conjecture.

\begin{conjecture}[Tachikawa {\cite[\S8]{Tachikawa}}]
\label{conj:tc2}
Let $\Lambda$ be a self-injective artin algebra, and let $X$ be a finitely
generated right $\Lambda$-module.  If
\[
 \Ext_\Lambda^i(X,X)=0\qquad\text{for every }i>0,
\]
then $X$ is projective.
\end{conjecture}

We say that an artin algebra $\Lambda$ satisfies
the \emph{Auslander--Reiten conjecture} \ARC{} if every finitely generated
right $\Lambda$-module $X$ satisfying \eqref{eq:arc} is projective, and,
when $\Lambda$ is self-injective, that $\Lambda$ satisfies \emph{Tachikawa's
second conjecture} \TCtwo{} if every finitely generated right
$\Lambda$-module $X$ with $\Ext_\Lambda^i(X,X)=0$ for every $i>0$ is
projective.  Thus Conjectures~\ref{conj:arc} and~\ref{conj:tc2} assert that
every artin algebra satisfies \ARC{} and that every self-injective artin
algebra satisfies \TCtwo.  Our main result is the
following.

\begin{theoremA}[= Corollary~\ref{thm:main}]
Let $R$ be a commutative artinian ring.  If every self-injective artin
$R$-algebra satisfies \TCtwo, then every artin $R$-algebra satisfies \ARC.
\end{theoremA}

The proof uses the two-fold trivial extension
\[
 T_2(\Lambda)=\begin{pmatrix}\Lambda&D\Lambda\\D\Lambda&\Lambda\end{pmatrix},
 \qquad D\Lambda=\Hom_R(\Lambda,E),
\]
of an artin $R$-algebra $\Lambda$, where $E$ is an injective envelope of
$R/\rad R$.  This is the algebra of $2\times2$ matrices with these entries in
which the product of two off-diagonal entries is zero.  It is self-injective,
and we show that if $T_2(\Lambda)$ satisfies \TCtwo, then $\Lambda$ satisfies
\ARC{} (Theorem~\ref{thm:transfer}).

We also consider the following classical homological conjectures.  For an
artin algebra $\Lambda$, let
\begin{equation}\label{eq:injective}
 0\longrightarrow\Lambda_\Lambda\longrightarrow I^0
 \longrightarrow I^1\longrightarrow I^2\longrightarrow\cdots
\end{equation}
be the minimal injective resolution of the right regular module.  We say
that $\Lambda$ satisfies
\begin{itemize}
\item the \emph{generalized Nakayama conjecture} \GNC{} \cite{AR} if every
      indecomposable injective right $\Lambda$-module is a direct summand of
      $I^i$ for some $i\geq0$;
\item the \emph{Auslander--Gorenstein conjecture} \AGC{} \cite[\S5]{AR94} if
      the following holds: if $\pd_\Lambda I^i\leq i$ for every $i\geq0$,
      then $\Lambda_\Lambda$ and ${}_\Lambda\Lambda$ have finite injective
      dimension;
\item the \emph{Nakayama conjecture} \NC{} \cite{Nakayama} if the following
      holds: if $I^i$ is projective for every $i\geq0$, then $\Lambda$ is
      self-injective;
\item the \emph{Gorenstein-projective conjecture} \GPC{} \cite{LH} if every
      finitely generated Gorenstein-projective right $\Lambda$-module $X$
      with $\Ext_\Lambda^i(X,X)=0$ for every $i>0$ is projective.
\end{itemize}
Each of these conjectures asserts that every artin algebra satisfies the
corresponding condition.

Theorem~A completes the following cycle of implications, where
$\Rightarrow$ denotes an implication for each individual algebra (each
self-injective algebra when \TCtwo{} is involved), and $\Rrightarrow$ denotes an
implication between the assertions that every artin $R$-algebra (every
self-injective artin $R$-algebra in the case of \TCtwo) satisfies the given
condition.  The implications other than Theorem~A are known and are
recalled in Proposition~\ref{prop:classical}.

\begin{samepage}
\begin{equation}\label{eq:cycle}
\xymatrix@C=3pc@R=2.5pc{
 \text{\GNC} \ar@{=>}[r] & \text{\AGC} \ar@{=>}[r] & \text{\NC} \ar@3{->}[d] \\
 \text{\ARC} \ar@3{<->}[u] \ar@{=>}[r] & \text{\GPC} \ar@{=>}[r]
 & \text{\TCtwo} \ar@3{->}@/^2pc/[ll]^{\text{Theorem A}}
}
\end{equation}

\begin{corollaryB}[= Corollary~\ref{cor:equivalences}]
Let $R$ be a commutative artinian ring.  The following statements are
equivalent.
\begin{enumerate}
\item Every self-injective artin $R$-algebra satisfies \TCtwo.
\item Every artin $R$-algebra satisfies \ARC.
\item Every artin $R$-algebra satisfies \GNC.
\item Every artin $R$-algebra satisfies \AGC.
\item Every artin $R$-algebra satisfies \NC.
\item Every artin $R$-algebra satisfies \GPC.
\end{enumerate}
\end{corollaryB}
\end{samepage}

We say that an artin algebra $\Lambda$ satisfies \emph{Tachikawa's first
conjecture} \TCone{} \cite[\S8]{Tachikawa} if the following holds: if
$\Ext_\Lambda^i(D\Lambda,\Lambda)=0$ for every $i>0$, then $\Lambda$ is
self-injective; the conjecture asserts that every artin algebra satisfies
\TCone.  An algebra satisfying \ARC{} also satisfies \TCone, and so
Theorem~A shows that Tachikawa's second conjecture implies his first: if
every self-injective artin $R$-algebra satisfies \TCtwo, then every artin
$R$-algebra satisfies \TCone{} (Corollary~\ref{cor:tc1}).

\begin{organization}
In Section~\ref{sec:induction}, Lemma~\ref{lem:induction} shows that
induction along an idempotent reflects projectivity, and
Lemma~\ref{lem:self-injective} shows that $T_2(\Lambda)$ is self-injective.
The main output of that section is Proposition~\ref{prop:transport}, which
identifies Ext groups over $T_2(\Lambda)$ of induced modules with Ext groups
over $\Lambda$ for modules $X$ with $\Ext_\Lambda^i(X,\Lambda)=0$ for every
$i>0$.  Section~\ref{sec:consequences} proves Theorem~\ref{thm:transfer} and
Theorem~A, recalls the known
implications, and deduces Corollary~B and the consequence for \TCone.
\end{organization}

\begin{conventions}
Throughout, $R$ is a commutative artinian ring and $E$ is an injective
envelope of $R/\rad R$.  An \emph{artin $R$-algebra} is an associative unital
$R$-algebra which is finitely generated as an $R$-module.  All algebras are
artin $R$-algebras, and all modules are finitely generated right modules
unless otherwise specified.  We write $\modu\Lambda$ for the category of
finitely generated right $\Lambda$-modules.  We write $D=\Hom_R(-,E)$; it is an exact duality between finitely
generated right and left $\Lambda$-modules, with $DDM\cong M$ naturally.  We give $D\Lambda$ the $\Lambda$-bimodule
structure
\[
 (a\varphi b)(c)=\varphi(bca)\qquad(a,b,c\in\Lambda,\ \varphi\in D\Lambda).
\]
\end{conventions}

\begin{aiuse}
The proof in this paper was found by GPT-6-Astra in research directed by the
author, and the first draft of the manuscript was written by GPT-6-Astra.
Claude Opus~5 revised the exposition of the manuscript.  The author set the
structure and much of the wording of the present text.  The author is
responsible for the result.
\end{aiuse}

\section{Induction to the two-fold trivial extension}
\label{sec:induction}

The proof of Theorem~A compares modules over $\Lambda$ with modules over an
algebra $\Gamma$ containing an idempotent $e$ with $e\Gamma e=\Lambda$, by
means of the functor $F$ in \eqref{eq:functors} below.
Lemma~\ref{lem:induction} gives the properties of $F$ used in the proof, for
an arbitrary idempotent; its Ext comparison assumes a Tor vanishing
condition.  The rest of the section treats the case $\Gamma=T_2(\Lambda)$, for
which Proposition~\ref{prop:transport} derives that condition from
$\Ext_\Lambda^i(X,\Lambda)=0$ for $i>0$.

Let $\Gamma$ be an algebra, let $e\in\Gamma$ be an idempotent, and put
$\Lambda=e\Gamma e$, an algebra with identity $e$.  Then $e\Gamma$ is a
$\Lambda$-$\Gamma$-bimodule, and we consider the functors
\begin{equation}\label{eq:functors}
 F=-\otimes_\Lambda e\Gamma\colon\modu\Lambda\longrightarrow\modu\Gamma,
 \qquad
 H=(-)e\colon\modu\Gamma\longrightarrow\modu\Lambda.
\end{equation}
We call $F$ \emph{induction along $e$}.  Since $\Hom_\Gamma(e\Gamma,Z)\cong
Ze$ for every $Z\in\modu\Gamma$, tensor--Hom adjunction gives isomorphisms
\begin{equation}\label{eq:adjunction}
 \Hom_\Gamma(FY,Z)\cong\Hom_\Lambda(Y,Ze)
\end{equation}
natural in $Y\in\modu\Lambda$ and $Z\in\modu\Gamma$, so $F$ is left adjoint
to $H$.  For every $Y\in\modu\Lambda$, since $e\Gamma e=\Lambda$,
\[
 (FY)e=Y\otimes_\Lambda e\Gamma e=Y\otimes_\Lambda\Lambda\cong Y,
\]
and under this isomorphism the unit $Y\to(FY)e$, $y\mapsto y\otimes e$, of the
adjunction is the identity.  In particular, the unit is an isomorphism.

\begin{lemma}\label{lem:induction}
In the setting of \eqref{eq:functors}, the following hold.
\begin{enumerate}
\item The functor $F$ is fully faithful.  A module $Y\in\modu\Lambda$ is
      projective if and only if $FY$ is projective.
\item Let $X\in\modu\Lambda$ satisfy $\Tor_i^\Lambda(X,e\Gamma)=0$ for every
      $i>0$.  Then for every $Z\in\modu\Gamma$ and every $n\geq0$, there is
      an isomorphism
      \[
       \Ext_\Gamma^n(FX,Z)\cong\Ext_\Lambda^n(X,Ze).
      \]
\end{enumerate}
\end{lemma}

\begin{proof}
(1): A left adjoint is fully faithful if and only if the unit of the
adjunction is an isomorphism, so $F$ is fully faithful.  Since
$F\Lambda=e\Gamma$ is a direct summand of $\Gamma_\Gamma$, the functor $F$
preserves projective modules.  Conversely, suppose that $FY$ is projective.
Choose an epimorphism $\pi\colon\Lambda^m\to Y$.  Since $F$ is right exact,
$F\pi\colon(e\Gamma)^m\to FY$ is an epimorphism, so there is $s\colon
FY\to(e\Gamma)^m$ with $F\pi\circ s=1_{FY}$.  Since $F$ is fully faithful,
$s=Fs'$ for some $s'\colon Y\to\Lambda^m$, and $F(\pi\circ s')=F\pi\circ
s=1_{FY}$ gives $\pi\circ s'=1_Y$.  Hence $Y$ is a direct
summand of $\Lambda^m$, and so $Y$ is projective.

(2): Choose a projective resolution $P_\bullet\to X$ in $\modu\Lambda$.  The
homology of $P_\bullet\otimes_\Lambda e\Gamma$ is $FX$ in degree $0$ and
$\Tor_i^\Lambda(X,e\Gamma)=0$ in degree $i>0$.  Hence
\[
 \cdots\longrightarrow FP_2\longrightarrow FP_1
 \longrightarrow FP_0\longrightarrow FX\longrightarrow0
\]
is exact.  Its terms $FP_j$ are projective by~(1), so $FP_\bullet$ is a
projective resolution of $FX$.  The isomorphisms \eqref{eq:adjunction} are
natural in $Y$, so they give an isomorphism of cochain complexes
\[
 \Hom_\Gamma(FP_\bullet,Z)\cong\Hom_\Lambda(P_\bullet,Ze).
\]
Taking cohomology in degree $n$ gives
$\Ext_\Gamma^n(FX,Z)\cong\Ext_\Lambda^n(X,Ze)$.
\end{proof}

We now construct, for an arbitrary algebra $\Lambda$, a self-injective
algebra $\Gamma$ and an idempotent $e$ with $e\Gamma e=\Lambda$.

\begin{definition}\label{def:two-fold}
Let $\Lambda$ be an artin $R$-algebra.  The \emph{two-fold trivial extension}
of $\Lambda$ is the $R$-module
\begin{equation}\label{eq:extension}
 T_2(\Lambda)=
 \begin{pmatrix}\Lambda&D\Lambda\\D\Lambda&\Lambda\end{pmatrix}
\end{equation}
with multiplication
\begin{equation}\label{eq:multiplication}
 \begin{pmatrix}a&\varphi\\\psi&b\end{pmatrix}
 \begin{pmatrix}a'&\varphi'\\\psi'&b'\end{pmatrix}
 =
 \begin{pmatrix}
 aa'&a\varphi'+\varphi b'\\
 \psi a'+b\psi'&bb'
 \end{pmatrix},
\end{equation}
where the products in the off-diagonal entries are the bimodule actions on
$D\Lambda$.
\end{definition}

For finite-dimensional algebras over a field, $T_2(\Lambda)$ is the case
$r=2$ of the $r$-fold trivial extension; see, for example,
\cite[\S2.2]{CDIM}.

The off-diagonal matrices form a submodule $I$, which is a
$(\Lambda\times\Lambda)$-bimodule via
\[
 (a,b)\begin{pmatrix}0&\varphi\\\psi&0\end{pmatrix}(a',b')
 =\begin{pmatrix}0&a\varphi b'\\b\psi a'&0\end{pmatrix}.
\]
By \eqref{eq:multiplication}, $T_2(\Lambda)$ is the trivial extension of
$\Lambda\times\Lambda$ by $I$, in which products of two elements of $I$ are
zero.  Hence $T_2(\Lambda)$ is an associative $R$-algebra with identity
$\operatorname{diag}(1,1)$, and it is finitely generated as an $R$-module
since $\Lambda$ and $D\Lambda$ are.  For $e=\operatorname{diag}(1,0)$, we have
$eT_2(\Lambda)e=\Lambda$, and
\begin{equation}\label{eq:eGamma}
 eT_2(\Lambda)=\begin{pmatrix}\Lambda&D\Lambda\end{pmatrix}
 \cong\Lambda\oplus D\Lambda
 \qquad\text{as left }\Lambda\text{-modules.}
\end{equation}

\begin{lemma}\label{lem:self-injective}
The artin $R$-algebra $T_2(\Lambda)$ is self-injective.
\end{lemma}

\begin{proof}
Put $\Gamma=T_2(\Lambda)$.  We show that
$\Gamma_\Gamma\cong D({}_\Gamma\Gamma)$.  Define an $R$-linear map
$\ell\colon\Gamma\to E$ by
\[
 \ell\!\begin{pmatrix}a&\varphi\\\psi&b\end{pmatrix}
 =\varphi(1)+\psi(1).
\]
For
$x=\begin{pmatrix}a&\varphi\\\psi&b\end{pmatrix}$ and
$y=\begin{pmatrix}a'&\varphi'\\\psi'&b'\end{pmatrix}$ in $\Gamma$, the
multiplication \eqref{eq:multiplication} and the bimodule structure on
$D\Lambda$ give
\[
 \ell(xy)=\varphi'(a)+\varphi(b')+\psi(a')+\psi'(b).
\]
Suppose that $\ell(xy)=0$ for every $y\in\Gamma$.  Taking $y$ with only
$\varphi'$, only $b'$, only $a'$, or only $\psi'$ nonzero shows that
$\varphi'(a)=0$ for every $\varphi'\in D\Lambda$, $\varphi=0$, $\psi=0$, and
$\psi'(b)=0$ for every $\psi'\in D\Lambda$.  Since $E$ is an injective
cogenerator of $R$-modules, the first and last conditions give $a=0$ and
$b=0$.  Hence the right $\Gamma$-module homomorphism
\[
 \Gamma_\Gamma\longrightarrow D({}_\Gamma\Gamma),\qquad
 x\longmapsto\bigl(y\longmapsto\ell(xy)\bigr)
\]
is injective, where $(h r)(y)=h(ry)$ for $h\in D\Gamma$ and $r,y\in\Gamma$.
Since $D$ preserves the length of $R$-modules, both sides have the same
length as $R$-modules, so this map is an isomorphism.  The right $\Gamma$-module $D({}_\Gamma\Gamma)$ is injective,
since $D$ sends the projective left module ${}_\Gamma\Gamma$ to an injective
right module.  Therefore $\Gamma_\Gamma$ is injective.
\end{proof}

For $\Gamma=T_2(\Lambda)$, the next proposition shows that the hypothesis
$\Ext_\Lambda^i(X,\Lambda)=0$ for $i>0$ implies the Tor vanishing of
Lemma~\ref{lem:induction}(2), and applies that lemma with $Z=FY$.

\begin{proposition}\label{prop:transport}
Let $\Gamma=T_2(\Lambda)$ and $e=\operatorname{diag}(1,0)$, and let $F$ be
as in \eqref{eq:functors}.  Let $X\in\modu\Lambda$ satisfy
$\Ext_\Lambda^i(X,\Lambda)=0$ for every $i>0$.  Then for every
$Y\in\modu\Lambda$ and every $n\geq0$,
\[
 \Ext_\Gamma^n(FX,FY)\cong\Ext_\Lambda^n(X,Y).
\]
\end{proposition}

\begin{proof}
Choose a projective resolution $P_\bullet\to X$ in $\modu\Lambda$ by finitely
generated projective modules.  Tensor--Hom adjunction and the isomorphism
$DD\Lambda\cong\Lambda$ of right $\Lambda$-modules give an isomorphism of
cochain complexes
\[
 D(P_\bullet\otimes_\Lambda D\Lambda)
 \cong\Hom_\Lambda(P_\bullet,DD\Lambda)
 \cong\Hom_\Lambda(P_\bullet,\Lambda).
\]
Since $D$ is exact, taking cohomology gives
\[
 D\Tor_i^\Lambda(X,D\Lambda)\cong\Ext_\Lambda^i(X,\Lambda)=0
 \qquad(i>0),
\]
and hence $\Tor_i^\Lambda(X,D\Lambda)=0$ for $i>0$.  As
$\Tor_i^\Lambda(X,\Lambda)=0$ for $i>0$, the decomposition \eqref{eq:eGamma}
gives $\Tor_i^\Lambda(X,e\Gamma)=0$ for every $i>0$.
Lemma~\ref{lem:induction}(2) with $Z=FY$ gives
$\Ext_\Gamma^n(FX,FY)\cong\Ext_\Lambda^n(X,(FY)e)$, and $(FY)e\cong Y$ by the
unit isomorphism stated before Lemma~\ref{lem:induction}.
\end{proof}

\section{Proof of Theorem A and the equivalence of the conjectures}
\label{sec:consequences}

This section first proves, for an individual algebra, that \TCtwo{} for its
two-fold trivial extension implies \ARC, and deduces Theorem~A.  It then
recalls the known implications among the conditions defined in the
Introduction, deduces Corollary~B, and proves the consequence for \TCone.

\begin{theorem}\label{thm:transfer}
Let $\Lambda$ be an artin algebra.  If $T_2(\Lambda)$ satisfies \TCtwo, then
$\Lambda$ satisfies \ARC.
\end{theorem}

\begin{proof}
Let $X\in\modu\Lambda$ satisfy \eqref{eq:arc}.  Put $\Gamma=T_2(\Lambda)$
and $e=\operatorname{diag}(1,0)$, and let $F=-\otimes_\Lambda e\Gamma$.  By
Lemma~\ref{lem:self-injective}, $\Gamma$ is self-injective.  Since $\Ext_\Lambda^i(X,\Lambda)=0$ for $i>0$,
Proposition~\ref{prop:transport} with $Y=X$ gives
\[
 \Ext_\Gamma^i(FX,FX)\cong\Ext_\Lambda^i(X,X)=0\qquad(i>0).
\]
Since $\Gamma$ satisfies \TCtwo, the module $FX$ is projective.  By
Lemma~\ref{lem:induction}(1), $X$ is projective.
\end{proof}

\begin{corollary}\label{thm:main}
Let $R$ be a commutative artinian ring.  If every self-injective artin
$R$-algebra satisfies \TCtwo, then every artin $R$-algebra satisfies \ARC.
\end{corollary}

\begin{proof}
Let $\Lambda$ be an artin $R$-algebra.  By Lemma~\ref{lem:self-injective},
$T_2(\Lambda)$ is a self-injective artin $R$-algebra, so it satisfies \TCtwo{} by
hypothesis.  Theorem~\ref{thm:transfer} shows that $\Lambda$ satisfies \ARC.
\end{proof}

Theorem~\ref{thm:transfer} passes from $\Lambda$ to the different algebra
$T_2(\Lambda)$, which is why the arrow labelled Theorem~A in \eqref{eq:cycle} is an
implication between assertions about all algebras.

To deduce Corollary~B, we recall the known implications in
\eqref{eq:cycle}.

\begin{proposition}\label{prop:classical}
Let $R$ be a commutative artinian ring.
\begin{enumerate}
\item Every artin $R$-algebra satisfies \ARC{} if and only if every artin
      $R$-algebra satisfies \GNC.
\item Let $\Lambda$ be an artin $R$-algebra.  If $\Lambda$ satisfies \GNC,
      then $\Lambda$ satisfies \AGC.  If $\Lambda$ satisfies \AGC, then
      $\Lambda$ satisfies \NC.
\item Suppose that every artin $R$-algebra satisfies \NC.  Then every
      self-injective artin $R$-algebra satisfies \TCtwo.
\item Let $\Lambda$ be an artin $R$-algebra.  If $\Lambda$ satisfies \ARC,
      then $\Lambda$ satisfies \GPC.
\end{enumerate}
\end{proposition}

\begin{proof}
(1): See \cite[Theorem~1.1]{AR}, which states this with \ARC{} replaced by the
condition that every generator $M$ with $\Ext^i(M,M)=0$ for every $i>0$ is
projective.  The two conditions agree for each algebra
\cite[pp.~283--284]{ADS}.

(2): See \cite[\S3.3]{TH}.

(3): By \cite[Corollary~1.2 and the remark following it]{AR}, the hypothesis
implies that every generator-cogenerator $M$ over an artin $R$-algebra with
$\Ext^i(M,M)=0$ for every $i>0$ is projective.  If $\Lambda$ is
self-injective and $\Ext_\Lambda^i(X,X)=0$ for every $i>0$, apply this to
$M=\Lambda\oplus X$.  It is a generator, and it is a cogenerator because
every indecomposable injective right module over the self-injective algebra
$\Lambda$ is projective and hence a direct summand of $\Lambda_\Lambda$.
Since $\Lambda_\Lambda$ is projective and injective,
$\Ext_\Lambda^i(M,M)\cong\Ext_\Lambda^i(X,X)=0$ for $i>0$.  Hence $M$ is
projective, and so is its direct summand $X$.

(4): This is noted in \cite{LH}.  Every finitely generated
Gorenstein-projective module $X$ satisfies $\Ext_\Lambda^i(X,\Lambda)=0$ for
$i>0$.  If moreover $\Ext_\Lambda^i(X,X)=0$ for $i>0$, then $X$ satisfies
\eqref{eq:arc}.
\end{proof}

\begin{corollary}\label{cor:equivalences}
Let $R$ be a commutative artinian ring.  The following statements are
equivalent.
\begin{enumerate}
\item Every self-injective artin $R$-algebra satisfies \TCtwo.
\item Every artin $R$-algebra satisfies \ARC.
\item Every artin $R$-algebra satisfies \GNC.
\item Every artin $R$-algebra satisfies \AGC.
\item Every artin $R$-algebra satisfies \NC.
\item Every artin $R$-algebra satisfies \GPC.
\end{enumerate}
\end{corollary}

\begin{proof}
Corollary~\ref{thm:main} gives (1) $\Rightarrow$ (2).
Proposition~\ref{prop:classical}(1) gives (2) $\Rightarrow$ (3),
Proposition~\ref{prop:classical}(2) applied to each algebra gives
(3) $\Rightarrow$ (4) $\Rightarrow$ (5), and
Proposition~\ref{prop:classical}(3) gives (5) $\Rightarrow$ (1).
Proposition~\ref{prop:classical}(4) applied to each algebra gives
(2) $\Rightarrow$ (6).  Finally, (6) $\Rightarrow$ (1) holds because a self-injective artin algebra
satisfying \GPC{} satisfies \TCtwo: every finitely generated module over it
is Gorenstein-projective.
\end{proof}

Finally, Corollary~\ref{thm:main} shows that Tachikawa's second conjecture
implies his first.

\begin{corollary}\label{cor:tc1}
Let $R$ be a commutative artinian ring.  If every self-injective artin
$R$-algebra satisfies \TCtwo, then every artin $R$-algebra satisfies \TCone.
\end{corollary}

\begin{proof}
Let $\Lambda$ be an artin $R$-algebra with $\Ext_\Lambda^i(D\Lambda,\Lambda)=0$
for every $i>0$.  The right $\Lambda$-module $D\Lambda$ is injective, so
$X=D\Lambda$ satisfies \eqref{eq:arc}.  By Corollary~\ref{thm:main}, $\Lambda$
satisfies \ARC, so $D\Lambda$ is projective.  Every indecomposable injective
right $\Lambda$-module is a direct summand of $D\Lambda$, and is therefore
projective.  The numbers of isomorphism classes of indecomposable injective
and of indecomposable projective right $\Lambda$-modules are both equal to
the number of isomorphism classes of simple modules, so every indecomposable
projective right $\Lambda$-module is injective.  Hence $\Lambda$ is
self-injective.
\end{proof}

\begin{thebibliography}{99}

\bibitem{ADS}
M.~Auslander, S.~Ding, and {\O}.~Solberg,
\emph{Liftings and weak liftings of modules},
J. Algebra \textbf{156} (1993), no.~2, 273--317.

\bibitem{AR}
M.~Auslander and I.~Reiten,
\emph{On a generalized version of the Nakayama conjecture},
Proc. Amer. Math. Soc. \textbf{52} (1975), 69--74.

\bibitem{AR94}
M.~Auslander and I.~Reiten,
\emph{$k$-Gorenstein algebras and syzygy modules},
J. Pure Appl. Algebra \textbf{92} (1994), no.~1, 1--27.

\bibitem{CDIM}
A.~Chan, E.~Darp\"o, O.~Iyama, and R.~Marczinzik,
\emph{Periodic trivial extension algebras and fractionally Calabi--Yau
algebras}, Ann. Sci. \'Ec. Norm. Sup\'er. (4) \textbf{58} (2025), no.~2,
463--510.

\bibitem{LH}
R.~Luo and Z.~Huang,
\emph{When are torsionless modules projective?},
J. Algebra \textbf{320} (2008), no.~5, 2156--2164.

\bibitem{Nakayama}
T.~Nakayama,
\emph{On algebras with complete homology},
Abh. Math. Sem. Univ. Hamburg \textbf{22} (1958), 300--307.

\bibitem{Tachikawa}
H.~Tachikawa,
\emph{Quasi-Frobenius rings and generalizations: QF-3 and QF-1 rings},
Lecture Notes in Mathematics, vol.~351, Springer-Verlag, Berlin, 1973.
Notes by C.~M.~Ringel.

\bibitem{TH}
X.~Tang and Z.~Huang,
\emph{Higher differential objects in additive categories},
J. Algebra \textbf{549} (2020), 128--164.

\end{thebibliography}

\end{document}
